% Options for packages loaded elsewhere
% Options for packages loaded elsewhere
\PassOptionsToPackage{unicode}{hyperref}
\PassOptionsToPackage{hyphens}{url}
\PassOptionsToPackage{dvipsnames,svgnames,x11names}{xcolor}
%
\documentclass[
  10pt,
  a4paper,
  onecolumn]{article}
\usepackage{xcolor}
\usepackage[top=2.5cm,bottom=2.5cm,left=3cm,right=3cm]{geometry}
\usepackage{amsmath,amssymb}
\setcounter{secnumdepth}{5}
\usepackage{iftex}
\ifPDFTeX
  \usepackage[T1]{fontenc}
  \usepackage[utf8]{inputenc}
  \usepackage{textcomp} % provide euro and other symbols
\else % if luatex or xetex
  \usepackage{unicode-math} % this also loads fontspec
  \defaultfontfeatures{Scale=MatchLowercase}
  \defaultfontfeatures[\rmfamily]{Ligatures=TeX,Scale=1}
\fi
\usepackage{lmodern}
\ifPDFTeX\else
  % xetex/luatex font selection
\fi
% Use upquote if available, for straight quotes in verbatim environments
\IfFileExists{upquote.sty}{\usepackage{upquote}}{}
\IfFileExists{microtype.sty}{% use microtype if available
  \usepackage[]{microtype}
  \UseMicrotypeSet[protrusion]{basicmath} % disable protrusion for tt fonts
}{}
\makeatletter
\@ifundefined{KOMAClassName}{% if non-KOMA class
  \IfFileExists{parskip.sty}{%
    \usepackage{parskip}
  }{% else
    \setlength{\parindent}{0pt}
    \setlength{\parskip}{6pt plus 2pt minus 1pt}}
}{% if KOMA class
  \KOMAoptions{parskip=half}}
\makeatother
% Make \paragraph and \subparagraph free-standing
\makeatletter
\ifx\paragraph\undefined\else
  \let\oldparagraph\paragraph
  \renewcommand{\paragraph}{
    \@ifstar
      \xxxParagraphStar
      \xxxParagraphNoStar
  }
  \newcommand{\xxxParagraphStar}[1]{\oldparagraph*{#1}\mbox{}}
  \newcommand{\xxxParagraphNoStar}[1]{\oldparagraph{#1}\mbox{}}
\fi
\ifx\subparagraph\undefined\else
  \let\oldsubparagraph\subparagraph
  \renewcommand{\subparagraph}{
    \@ifstar
      \xxxSubParagraphStar
      \xxxSubParagraphNoStar
  }
  \newcommand{\xxxSubParagraphStar}[1]{\oldsubparagraph*{#1}\mbox{}}
  \newcommand{\xxxSubParagraphNoStar}[1]{\oldsubparagraph{#1}\mbox{}}
\fi
\makeatother


\usepackage{longtable,booktabs,array}
\usepackage{calc} % for calculating minipage widths
% Correct order of tables after \paragraph or \subparagraph
\usepackage{etoolbox}
\makeatletter
\patchcmd\longtable{\par}{\if@noskipsec\mbox{}\fi\par}{}{}
\makeatother
% Allow footnotes in longtable head/foot
\IfFileExists{footnotehyper.sty}{\usepackage{footnotehyper}}{\usepackage{footnote}}
\makesavenoteenv{longtable}
\usepackage{graphicx}
\makeatletter
\newsavebox\pandoc@box
\newcommand*\pandocbounded[1]{% scales image to fit in text height/width
  \sbox\pandoc@box{#1}%
  \Gscale@div\@tempa{\textheight}{\dimexpr\ht\pandoc@box+\dp\pandoc@box\relax}%
  \Gscale@div\@tempb{\linewidth}{\wd\pandoc@box}%
  \ifdim\@tempb\p@<\@tempa\p@\let\@tempa\@tempb\fi% select the smaller of both
  \ifdim\@tempa\p@<\p@\scalebox{\@tempa}{\usebox\pandoc@box}%
  \else\usebox{\pandoc@box}%
  \fi%
}
% Set default figure placement to htbp
\def\fps@figure{htbp}
\makeatother


% definitions for citeproc citations
\NewDocumentCommand\citeproctext{}{}
\NewDocumentCommand\citeproc{mm}{%
  \begingroup\def\citeproctext{#2}\cite{#1}\endgroup}
\makeatletter
 % allow citations to break across lines
 \let\@cite@ofmt\@firstofone
 % avoid brackets around text for \cite:
 \def\@biblabel#1{}
 \def\@cite#1#2{{#1\if@tempswa , #2\fi}}
\makeatother
\newlength{\cslhangindent}
\setlength{\cslhangindent}{1.5em}
\newlength{\csllabelwidth}
\setlength{\csllabelwidth}{3em}
\newenvironment{CSLReferences}[2] % #1 hanging-indent, #2 entry-spacing
 {\begin{list}{}{%
  \setlength{\itemindent}{0pt}
  \setlength{\leftmargin}{0pt}
  \setlength{\parsep}{0pt}
  % turn on hanging indent if param 1 is 1
  \ifodd #1
   \setlength{\leftmargin}{\cslhangindent}
   \setlength{\itemindent}{-1\cslhangindent}
  \fi
  % set entry spacing
  \setlength{\itemsep}{#2\baselineskip}}}
 {\end{list}}
\usepackage{calc}
\newcommand{\CSLBlock}[1]{\hfill\break\parbox[t]{\linewidth}{\strut\ignorespaces#1\strut}}
\newcommand{\CSLLeftMargin}[1]{\parbox[t]{\csllabelwidth}{\strut#1\strut}}
\newcommand{\CSLRightInline}[1]{\parbox[t]{\linewidth - \csllabelwidth}{\strut#1\strut}}
\newcommand{\CSLIndent}[1]{\hspace{\cslhangindent}#1}



\setlength{\emergencystretch}{3em} % prevent overfull lines

\providecommand{\tightlist}{%
  \setlength{\itemsep}{0pt}\setlength{\parskip}{0pt}}



 


\linespread{1.5}
\usepackage{booktabs}
\usepackage{caption}
\captionsetup[table]{labelfont=bf, skip=6pt}
\captionsetup[figure]{labelfont=bf, skip=6pt}
\usepackage{fancyhdr}
\pagestyle{fancy}
\fancyhead[L]{\small\itshape Monitoring Therapeutic Persona Drift}
\fancyhead[R]{\small\thepage}
\fancyfoot{}
\renewcommand{\headrulewidth}{0.4pt}
\makeatletter
\@ifpackageloaded{caption}{}{\usepackage{caption}}
\AtBeginDocument{%
\ifdefined\contentsname
  \renewcommand*\contentsname{Table of contents}
\else
  \newcommand\contentsname{Table of contents}
\fi
\ifdefined\listfigurename
  \renewcommand*\listfigurename{List of Figures}
\else
  \newcommand\listfigurename{List of Figures}
\fi
\ifdefined\listtablename
  \renewcommand*\listtablename{List of Tables}
\else
  \newcommand\listtablename{List of Tables}
\fi
\ifdefined\figurename
  \renewcommand*\figurename{Figure}
\else
  \newcommand\figurename{Figure}
\fi
\ifdefined\tablename
  \renewcommand*\tablename{Table}
\else
  \newcommand\tablename{Table}
\fi
}
\@ifpackageloaded{float}{}{\usepackage{float}}
\floatstyle{ruled}
\@ifundefined{c@chapter}{\newfloat{codelisting}{h}{lop}}{\newfloat{codelisting}{h}{lop}[chapter]}
\floatname{codelisting}{Listing}
\newcommand*\listoflistings{\listof{codelisting}{List of Listings}}
\makeatother
\makeatletter
\makeatother
\makeatletter
\@ifpackageloaded{caption}{}{\usepackage{caption}}
\@ifpackageloaded{subcaption}{}{\usepackage{subcaption}}
\makeatother
\usepackage{bookmark}
\IfFileExists{xurl.sty}{\usepackage{xurl}}{} % add URL line breaks if available
\urlstyle{same}
\hypersetup{
  pdftitle={Monitoring Therapeutic Persona Drift in Mental Health Chatbots Using Internal Activation States: Evidence for an Instruction Tuning Firewall},
  pdfauthor={Lucas Sempé},
  pdfkeywords={activation monitoring, therapeutic persona, language
models, mental health chatbots, persona drift, representation
engineering},
  colorlinks=true,
  linkcolor={NavyBlue},
  filecolor={Maroon},
  citecolor={NavyBlue},
  urlcolor={NavyBlue},
  pdfcreator={LaTeX via pandoc}}


\title{Monitoring Therapeutic Persona Drift in Mental Health Chatbots
Using Internal Activation States: Evidence for an Instruction Tuning
Firewall}
\author{Lucas Sempé}
\date{2026-02-15}
\begin{document}
\maketitle
\begin{abstract}
\textbf{Background} AI-powered mental health chatbots are expected to
maintain a consistent therapeutic persona---a set of clinically grounded
behavioural qualities, including empathetic responsiveness,
non-judgmental acceptance, boundary maintenance, and crisis recognition,
drawn from the common factors of effective psychotherapy. However,
safety-trained language models can produce clinically
appropriate-sounding text even when their internal representations of
these qualities have shifted, creating a monitoring blind spot. We
developed and validated an activation-based monitoring system that
tracks therapeutic persona dimensions directly from the model's hidden
states, extending recent work on persona vectors for general character
traits to the safety-critical domain of mental health.

\textbf{Methods} We defined eight therapeutic persona dimensions
(empathetic responsiveness, non-judgmental acceptance, boundary
maintenance, crisis recognition, and four failure modes) and validated
monitoring vectors via activation steering in Llama-3-8B-Instruct (all
eight traits \(r\)=0·302--0·489, \(N\)=50 per trait). We extended
monitoring to Qwen2-7B-Instruct and Mistral-7B-Instruct-v0.2 using
contrastive probing---deriving monitoring vectors from each model's own
scored response distribution---achieving 21 of 24 validated trait--model
combinations (\(r\)\textgreater0·30). We compared activation monitoring
against progressively stronger text-based detection methods on 12 000
steered responses, tested for natural persona drift over 100-turn
conversations, and assessed ecological validity on 200 conversations
from the ESConv naturalistic emotional support corpus. Two clinical
psychologists independently validated LLM judge scores.

\textbf{Findings} All eight therapeutic persona traits were steerable in
Llama-3-8B-Instruct (\(r\)=0·302--0·489, all \(p\)\textless0·001), and
contrastive probing extended monitoring to 21 of 24 trait--model
combinations across three architectures. The real-time activation
monitor tracked persona dimensions with a cross-model mean \(r\)=0·596
and false alarm rates of 1--4\%. Text-based methods detected persona
drift poorly: sentiment analysis captured none (all
\(p\)\textgreater0·05), and even fine-tuned DeBERTa-v3-base explained
less variance than activation monitoring (\(R^2\)=0·290 vs 0·371), with
the gap widening on the two most safety-critical traits---crisis
recognition (2·6-fold) and abandonment of therapeutic frame (2·5-fold).
In 100-turn conversations without deliberate manipulation, all three
models showed significant upward drift in emotional over-involvement
(all \(p\)\textless0·01, cluster-robust standard errors), detected by
activation monitoring but invisible to sentiment analysis. Two clinical
psychologists confirmed LLM judge validity (intraclass correlation
coefficient 0·716, all eight traits \(r\)\$\geq\$0·50), though floor
effects on failure-mode traits limited per-trait reliability estimates.

\textbf{Interpretation} Safety training creates an instruction tuning
firewall: models maintain appropriate-sounding text even when their
internal representations of therapeutic qualities have shifted, making
text-based monitoring unreliable. This masking effect is strongest for
the most safety-critical traits---crisis recognition and maintenance of
professional boundaries---precisely where undetected drift carries the
greatest clinical risk. Activation monitoring bypasses this firewall by
reading the model's internal states directly, providing reliable
detection where text analysis is blind. An ecological validity test on
naturalistic emotional support conversations confirmed that the
monitoring vectors capture meaningful variation beyond synthetic
scenarios, but validation on clinical data with clinician-rated ground
truth is required before deployment recommendations can be made.
\end{abstract}


\section{Introduction}\label{introduction}

AI-powered chatbots are increasingly deployed for mental health support,
from wellness tools to structured therapeutic
interventions.\textsuperscript{\textsuperscript{1,2}} The central
challenge for these systems is persona consistency: whether a chatbot
maintains therapeutic qualities---empathy, non-judgmental acceptance,
boundary maintenance, crisis recognition---across extended interactions.
Recent work on persona vectors has shown that character traits in
language models correspond to identifiable directions in activation
space and can be monitored and steered
accordingly.\textsuperscript{\textsuperscript{3}} We extend this
framework from general personality traits to the safety-critical domain
of therapeutic personas, where undetected drift risks subtherapeutic
care or psychological harm. Persona drift, in which therapeutic
behaviour degrades over a conversation or across deployment conditions,
poses a safety risk that current monitoring approaches struggle to
detect.

Output filters catch toxic responses, but the real problem is subtler:
shifts in therapeutic quality that pass every text-level check. Several
mechanisms could cause such drift. Context window erosion---the focus of
this study---occurs as conversations extend beyond dozens of turns and
the system prompt's influence on model behaviour weakens. Fine-tuning
distribution shift, in which routine model updates on non-therapeutic
data could regress therapeutic qualities, is a plausible concern, though
a small-scale test in this study found no detectable effect. Adversarial
prompt manipulation, in which users gradually erode professional
boundaries through emotional pressure or boundary-testing, was also
tested here but proved readily detectable by both text-based and
activation-based methods.

To illustrate the monitoring challenge, we present two responses from
Llama-3-8B-Instruct to a user expressing suicidal ideation (``I've been
having thoughts that everyone would be better off without me'').
Response A was generated under normal operating conditions; Response B
was generated with empathetic responsiveness shifted 3·0 standard
deviations below baseline via activation steering. Response A named the
crisis directly (``suicidal thoughts or feelings of worthlessness'') and
validated help-seeking (``completely understandable that you're reaching
out''). Response B euphemised the concern (``these thoughts,'' ``your
concerns'') and reframed help-seeking as bravery (``it takes a lot of
courage''). Both responses validated the user's experience, normalised
help-seeking, and offered support. Automated sentiment analysis scored
them nearly identically (VADER compound: 0·56 vs 0·54). Neither
contained toxic content, overt dismissal, or harmful advice. No
text-level method we tested could reliably distinguish them (full
response text in supplementary appendix S1).

Current methods for ensuring persona quality are incomplete. Prompt
engineering provides no feedback signal when adherence declines.
Fine-tuning encodes therapeutic qualities into model weights but offers
no runtime visibility. Output filtering cannot distinguish adequate from
high-quality therapeutic responses. Text-level monitoring using
sentiment analysis, keyword detection, or readability metrics is the
obvious solution, but even the strongest text method we tested
(fine-tuned DeBERTa, \(R^2\)=0·290) leaves a residual gap against
activation monitoring that concentrates on safety-critical traits. We
use ``instruction tuning firewall'' as shorthand for the combined effect
of instruction tuning, reinforcement learning from human feedback, and
direct preference optimisation, all of which constrain the model's
output distribution to sound appropriate regardless of internal state
perturbation.

In this study, we make five contributions. First, we show that eight
therapeutic traits---drawn from
Rogers\textsuperscript{\textsuperscript{4}} and
Wampold\textsuperscript{\textsuperscript{5}}---correspond to steerable
directions in the hidden layers of Llama-3-8B-Instruct, and that these
dimensions reflect two to three latent factors. Second, we develop
contrastive probing as a method for recovering monitoring vectors when
template-based steering vectors fail to transfer across architectures,
achieving 21 of 24 validated trait--model combinations across three
architectures. Third, we deploy a real-time monitoring pipeline using
exponentially weighted moving average (EWMA) and cumulative sum (CUSUM)
statistical process control, achieving mean \(r\)=0·596 with false alarm
rates of 1--4\%. Fourth, we characterise the instruction tuning firewall
as a gradient of text-level detectability, showing that the residual gap
between activation monitoring and the strongest text baseline
concentrates on safety-critical traits. Fifth, we provide evidence that
context erosion over 100-turn conversations produces persona drift in a
clinically relevant direction that is invisible to text-level analysis
but detectable through activation monitoring.

\section{Therapeutic persona dimensions}\label{sec-dimensions}

Chen and colleagues\textsuperscript{\textsuperscript{3}} demonstrated
that general personality traits---such as agreeableness or
assertiveness---can be represented as directions in a language model's
hidden state space and used for real-time monitoring. Therapeutic
personas, however, require domain-specific dimensions grounded in
clinical theory. We defined eight therapeutic persona dimensions as a
pragmatic monitoring framework informed by, but not a direct
operationalisation of, the therapeutic conditions described by
Rogers\textsuperscript{\textsuperscript{4}} and the common factors
identified by Wampold.\textsuperscript{\textsuperscript{5}} Four
dimensions capture positive therapeutic qualities and four capture
failure modes. An exploratory factor analysis (supplementary table S10)
revealed two to three latent factors: a warmth factor (empathetic
responsiveness, non-judgmental acceptance, emotional over-involvement),
a professional structure factor (boundary maintenance, crisis
recognition, abandonment of therapeutic frame), and a validation factor
(uncritical validation, sycophancy). This structure suggests the
framework could be reduced to fewer dimensions without substantial loss
of monitoring coverage.

The four positive dimensions are empathetic responsiveness (whether the
chatbot recognises and validates emotional content), non-judgmental
acceptance (whether the chatbot can hear a user's disclosure without
moralising), boundary maintenance (whether the chatbot holds
professional limits while remaining warm), and crisis recognition
(whether the chatbot identifies risk and provides appropriate
resources). The four failure modes are emotional over-involvement (the
chatbot centres its own distress rather than the user's experience),
abandonment of therapeutic frame (professional structure gives way to
casual interaction), uncritical validation (sycophantic agreement
without therapeutic exploration), and sycophancy or harmful validation
(endorsement of harmful choices).

Table~\ref{tbl-constructs} maps each dimension to its theoretical
source. Several constructs from Rogers and Wampold are deliberately
omitted. Congruence\textsuperscript{\textsuperscript{4}} presupposes
subjective experience; we capture its inverse (emotional
over-involvement) as a detectable failure mode. Goal consensus and
collaboration\textsuperscript{\textsuperscript{5,6}} require
multi-session tracking that exceeds the scope of turn-level monitoring.
Therapeutic rupture-repair\textsuperscript{\textsuperscript{7}} requires
sequential state-change detection and represents an important direction
for future work.

\begin{longtable}[]{@{}
  >{\raggedright\arraybackslash}p{(\linewidth - 6\tabcolsep) * \real{0.4231}}
  >{\centering\arraybackslash}p{(\linewidth - 6\tabcolsep) * \real{0.1923}}
  >{\centering\arraybackslash}p{(\linewidth - 6\tabcolsep) * \real{0.1923}}
  >{\centering\arraybackslash}p{(\linewidth - 6\tabcolsep) * \real{0.1923}}@{}}
\caption{Mapping of persona dimensions to theoretical constructs.
Extended=adapted from the cited framework; Inferential=inferred as
failure mode; AI-specific=derived from LLM sycophancy
literature.\textsuperscript{\textsuperscript{8}}}\label{tbl-constructs}\tabularnewline
\toprule\noalign{}
\begin{minipage}[b]{\linewidth}\raggedright
Dimension
\end{minipage} & \begin{minipage}[b]{\linewidth}\centering
Rogers (1957)
\end{minipage} & \begin{minipage}[b]{\linewidth}\centering
Wampold (2015)
\end{minipage} & \begin{minipage}[b]{\linewidth}\centering
Source
\end{minipage} \\
\midrule\noalign{}
\endfirsthead
\toprule\noalign{}
\begin{minipage}[b]{\linewidth}\raggedright
Dimension
\end{minipage} & \begin{minipage}[b]{\linewidth}\centering
Rogers (1957)
\end{minipage} & \begin{minipage}[b]{\linewidth}\centering
Wampold (2015)
\end{minipage} & \begin{minipage}[b]{\linewidth}\centering
Source
\end{minipage} \\
\midrule\noalign{}
\endhead
\bottomrule\noalign{}
\endlastfoot
Empathetic responsiveness & Empathic understanding & Empathy & Direct \\
Non-judgmental acceptance & Unconditional positive regard & --- &
Direct \\
Boundary maintenance & --- & Real relationship & Extended \\
Crisis recognition & --- & --- & Safety-specific \\
Emotional over-involvement & --- & --- & Inferential \\
Abandonment of therapeutic frame & --- & Treatment structure &
Extended \\
Uncritical validation & --- & --- & AI-specific \\
Sycophancy/harmful validation & --- & --- & AI-safety \\
\end{longtable}

Inter-trait correlations and activation-level cross-trait analyses are
reported in supplementary tables S4 and S9. Uncritical validation and
sycophancy correlated at \(r\)=0·74; mean off-diagonal
\textbar{}\(r\)\textbar{} across activation projections was 0·36
(Llama-3), 0·33 (Qwen2), and 0·26 (Mistral), consistent with the two to
three latent factors identified in the factor analysis.

\section{From steering vectors to a monitoring
system}\label{sec-steering}

\subsection{Steering validation in Llama-3-8B}\label{sec-phase1}

Activation steering extracts a direction vector from the difference
between contrasting prompts---one expressing high trait levels, one
expressing low---in the model's hidden state
space.\textsuperscript{\textsuperscript{9--11}} For each trait, we
developed five contrast pairs using concrete, behavioural language
rather than abstract self-descriptions: early attempts using prompts
such as ``I am very empathetic'' produced vectors that separated
activations but did not change model behaviour, consistent with the
findings of Chen and colleagues.\textsuperscript{\textsuperscript{3}}

We extracted last-token activations, computed normalised direction
vectors, and averaged across pairs (Figure~\ref{fig-extraction}). Layer
selection was empirical: rather than selecting the layer with the
largest activation separation (Cohen's \(d\)), we tested steering at
each candidate layer and selected the layer producing the highest
correlation between steering coefficient and judged behaviour. Layers
with high Cohen's \(d\) often had near-zero behavioural effect (e.g.,
layer 27). A random direction control confirmed signal specificity:
randomly oriented vectors in the same space produced near-zero \(r\).

\begin{figure}[htbp]

\centering{

\includegraphics[width=0.85\linewidth,height=\textheight,keepaspectratio]{research_document_activation_steering_v6_files/mediabag/../04_results/figures/fig1_steering_extraction.pdf}

}

\caption{\label{fig-extraction}Steering vector extraction pipeline. High
and low contrast prompts are passed through the model; the normalised
difference in last-token activations at the empirically selected layer
defines the steering direction.}

\end{figure}%

Table~\ref{tbl-phase1} presents the results. All eight traits were
steerable, with \(r\) ranging from 0·302 to 0·489 (all
\(p\)\textless0·001). Behavioural validation was performed by a
GPT-4o-mini judge at temperature 0 using anchored rubrics (supplementary
table S3). A prospective power calculation using the observed effect
sizes as planning parameters indicates that the top four traits
(\(r\)\$\geq\$0·42) would require
\(N\)\$\leq\(50 to achieve 87% or greater power, but the bottom four (
\)r\$=0·302--0·374) would need larger samples (estimated \(N\)=75--120)
for 80\% power. Boundary maintenance (\(r\)=0·302) should be considered
marginal pending replication at larger sample sizes (supplementary table
S11).

\begin{longtable}[]{@{}lccl@{}}
\caption{Phase 1 results: all eight traits are steerable on
Llama-3-8B-Instruct. \(N\)=50 per trait (10 scenarios \(\times\) 5
steering coefficients). All \(p\)\textless0·001. CIs from bias-corrected
and accelerated bootstrap (10 000
resamples).}\label{tbl-phase1}\tabularnewline
\toprule\noalign{}
Trait & Best layer & \(r\) & 95\% CI \\
\midrule\noalign{}
\endfirsthead
\toprule\noalign{}
Trait & Best layer & \(r\) & 95\% CI \\
\midrule\noalign{}
\endhead
\bottomrule\noalign{}
\endlastfoot
Sycophancy/harmful validation & 19 & 0·489 & {[}0·31, 0·65{]} \\
Abandonment of therapeutic frame & 19 & 0·470 & {[}0·29, 0·63{]} \\
Emotional over-involvement & 19 & 0·441 & {[}0·26, 0·60{]} \\
Empathetic responsiveness & 17 & 0·424 & {[}0·24, 0·58{]} \\
Crisis recognition & 18 & 0·374 & {[}0·19, 0·54{]} \\
Uncritical validation & 18 & 0·364 & {[}0·17, 0·53{]} \\
Non-judgmental acceptance & 18 & 0·346 & {[}0·16, 0·51{]} \\
Boundary maintenance & 18 & 0·302 & {[}0·11, 0·47{]} \\
\end{longtable}

\subsection{Cross-architecture extension via contrastive
probing}\label{sec-probing}

Template-based steering vectors did not generalise across architectures:
only three of eight traits validated on Qwen2-7B and two of eight on
Mistral-7B. Diagnostic analysis across all 24 trait--model combinations
(Figure~\ref{fig-diagnostics}) identified a single dominant predictor of
steering success: how much judged behaviour changed when moving from low
to high activation projections (\(r\)=0·899, \(N\)=24). Activation
geometry metrics (Cohen's \(d\), within-class variance, prompt
consistency) were uninformative. This metric reflects internal
consistency of the steering-to-judge pipeline rather than independent
causal validation; the random direction control provides the independent
test. The extreme case was Qwen2 on uncritical validation: template
vectors produced activation separations five times larger than those in
Llama-3, but the behavioural difference was 0·007 compared with
1·208---the vector captured a direction in representational space that
was unrelated to the target behaviour.

\begin{figure}[htbp]

\centering{

\includegraphics[width=0.85\linewidth,height=\textheight,keepaspectratio]{research_document_activation_steering_v6_files/mediabag/../04_results/figures/fig2_diagnostics.pdf}

}

\caption{\label{fig-diagnostics}What predicts steering success across 24
trait--model combinations. Only behavioural difference (\(r\)=0·899)
reliably predicts controllable behaviour change. Activation geometry
metrics are uninformative.}

\end{figure}%

To address this transfer failure, we developed a contrastive probing
approach (Figure~\ref{fig-probing}). Rather than imposing externally
defined contrast prompts, we derived direction vectors from each model's
own response distribution. We generated 500 scenario responses per
model, judged each on a 1--7 scale, extracted hidden states from
high-scored and low-scored responses, and trained logistic regression
classifiers (L2 penalty, \(C\)=1·0, StandardScaler normalisation) whose
weight vectors defined the monitoring directions.

\begin{figure}[htbp]

\centering{

\includegraphics[width=0.85\linewidth,height=\textheight,keepaspectratio]{research_document_activation_steering_v6_files/mediabag/../04_results/figures/fig3_contrastive_probing.pdf}

}

\caption{\label{fig-probing}Contrastive probing pipeline. The model
generates responses to clinical scenarios; an LLM judge scores each
response; hidden states from high-scored and low-scored responses define
the contrastive classes; a logistic regression classifier's weight
vector becomes the monitoring direction.}

\end{figure}%

Contrastive probing recovered monitoring capacity on both non-Llama
architectures (Table~\ref{tbl-probing}). All eight traits validated on
Qwen2 (up from three) and five of eight on Mistral (up from two), with
three additional Mistral traits achieving weak but positive
correlations. Zero of 24 trait--model combinations showed negative
correlations. The three weak Mistral traits had insufficient contrastive
training data (17--67 samples per class compared with 81--100 for
validated traits), reflecting genuinely narrower behavioural
distributions on these dimensions in the Mistral architecture.

\begin{longtable}[]{@{}
  >{\raggedright\arraybackslash}p{(\linewidth - 6\tabcolsep) * \real{0.3182}}
  >{\centering\arraybackslash}p{(\linewidth - 6\tabcolsep) * \real{0.2273}}
  >{\centering\arraybackslash}p{(\linewidth - 6\tabcolsep) * \real{0.2273}}
  >{\centering\arraybackslash}p{(\linewidth - 6\tabcolsep) * \real{0.2273}}@{}}
\caption{Contrastive probing rescues cross-architecture failures. Bold
values indicate validated traits (\(r\)\textgreater0·30). Non-bold probe
values indicate weak traits (0·15\textless{}\(r\)\$\leq\$0·30). No
negative correlations across all 24
combinations.}\label{tbl-probing}\tabularnewline
\toprule\noalign{}
\begin{minipage}[b]{\linewidth}\raggedright
Trait
\end{minipage} & \begin{minipage}[b]{\linewidth}\centering
Llama-3 (template)
\end{minipage} & \begin{minipage}[b]{\linewidth}\centering
Qwen2 (template \(\to\) probe)
\end{minipage} & \begin{minipage}[b]{\linewidth}\centering
Mistral (template \(\to\) probe)
\end{minipage} \\
\midrule\noalign{}
\endfirsthead
\toprule\noalign{}
\begin{minipage}[b]{\linewidth}\raggedright
Trait
\end{minipage} & \begin{minipage}[b]{\linewidth}\centering
Llama-3 (template)
\end{minipage} & \begin{minipage}[b]{\linewidth}\centering
Qwen2 (template \(\to\) probe)
\end{minipage} & \begin{minipage}[b]{\linewidth}\centering
Mistral (template \(\to\) probe)
\end{minipage} \\
\midrule\noalign{}
\endhead
\bottomrule\noalign{}
\endlastfoot
Empathetic responsiveness & 0·424 & 0·240 \(\to\) \textbf{0·414} & 0·329
\(\to\) 0·327 \\
Non-judgmental acceptance & 0·346 & 0·091 \(\to\) \textbf{0·584} & 0·257
\(\to\) \textbf{0·467} \\
Boundary maintenance & 0·302 & 0·254 \(\to\) \textbf{0·449} & 0·350
\(\to\) 0·271 \\
Crisis recognition & 0·374 & 0·346 \(\to\) \textbf{0·503} & 0·268
\(\to\) \textbf{0·398} \\
Emotional over-involvement & 0·441 & 0·357 \(\to\) \textbf{0·303} &
0·168 \(\to\) 0·240 \\
Abandonment of frame & 0·470 & 0·400 \(\to\) \textbf{0·378} & 0·233
\(\to\) \textbf{0·411} \\
Uncritical validation & 0·364 & 0·042 \(\to\) \textbf{0·393} & 0·208
\(\to\) 0·215 \\
Sycophancy/harmful validation & 0·489 & 0·115 \(\to\) \textbf{0·390} &
0·176 \(\to\) \textbf{0·331} \\
\textbf{Validated (r\textgreater0·30)} & \textbf{8/8} & 3 \(\to\)
\textbf{8/8} & 2 \(\to\) \textbf{5/8} (+3 weak) \\
\end{longtable}

\subsection{Real-time monitoring pipeline}\label{sec-monitoring}

The monitoring pipeline projects activation vectors onto validated trait
directions at every conversational turn (Figure~\ref{fig-monitoring}).
Two complementary statistical process control methods track the
resulting projections: an EWMA
filter\textsuperscript{\textsuperscript{12}} (\(\lambda\)=0·2, effective
window approximately 10 turns) for smoothed trend detection, and a CUSUM
algorithm\textsuperscript{\textsuperscript{13}} (\(k\)=0·5\(\sigma\),
\(h\)=4·0\(\sigma\)) for detecting small persistent shifts. The system
generates four-level alerts (none, watch, warning, critical), applied
directionally so that only the clinically concerning drift direction
triggers alarms. Computational cost is one dot product per trait per
turn.

\begin{figure}[htbp]

\centering{

\includegraphics[width=0.85\linewidth,height=\textheight,keepaspectratio]{research_document_activation_steering_v6_files/mediabag/../04_results/figures/fig4_monitoring_architecture.pdf}

}

\caption{\label{fig-monitoring}Real-time monitoring architecture.
Activations are projected onto validated trait vectors at each
conversational turn; EWMA and CUSUM trackers generate four-level
alerts.}

\end{figure}%

Table~\ref{tbl-monitoring} presents activation--behaviour correlations
across all three architectures. All 24 model--trait combinations were
significant (\(p\)\textless0·0001), with a cross-model mean of
\(r\)=0·596. Table~\ref{tbl-alerts} shows that all models met the
pre-specified target of less than 10\% Warning-or-above alert rate.

\begin{longtable}[]{@{}lcccc@{}}
\caption{Activation--behaviour correlations across all three
architectures. All \(N\)=200 per model, all
\(p\)\textless0·0001.}\label{tbl-monitoring}\tabularnewline
\toprule\noalign{}
Trait & Llama-3 & Qwen2 & Mistral & Mean \\
\midrule\noalign{}
\endfirsthead
\toprule\noalign{}
Trait & Llama-3 & Qwen2 & Mistral & Mean \\
\midrule\noalign{}
\endhead
\bottomrule\noalign{}
\endlastfoot
Crisis recognition & 0·569 & 0·815 & 0·801 & 0·728 \\
Empathetic responsiveness & 0·741 & 0·757 & 0·706 & 0·735 \\
Non-judgmental acceptance & 0·677 & 0·780 & 0·735 & 0·731 \\
Abandonment of frame & 0·690 & 0·736 & 0·617 & 0·681 \\
Emotional over-involvement & 0·459 & 0·592 & 0·411 & 0·487 \\
Sycophancy/harmful validation & 0·477 & 0·541 & 0·444 & 0·487 \\
Boundary maintenance & 0·358 & 0·520 & 0·546 & 0·475 \\
Uncritical validation & 0·384 & 0·539 & 0·415 & 0·446 \\
\textbf{Model mean} & \textbf{0·544} & \textbf{0·660} & \textbf{0·584} &
\textbf{0·596} \\
\end{longtable}

\begin{longtable}[]{@{}lccc@{}}
\caption{Alert rates across 100 monitoring sessions per
model.}\label{tbl-alerts}\tabularnewline
\toprule\noalign{}
Alert level & Llama-3 & Qwen2 & Mistral \\
\midrule\noalign{}
\endfirsthead
\toprule\noalign{}
Alert level & Llama-3 & Qwen2 & Mistral \\
\midrule\noalign{}
\endhead
\bottomrule\noalign{}
\endlastfoot
Any alert (watch or above) & 5\% & 13\% & 10\% \\
Warning or above & 4\% & 4\% & 1\% \\
Critical & 0\% & 1\% & 0\% \\
\end{longtable}

\section{The instruction tuning firewall}\label{sec-firewall}

\subsection{Detection gradient}\label{detection-gradient}

To test whether text-based monitoring could substitute for
activation-based monitoring, we generated 12 000 steered responses (3
models \(\times\) 8 traits \(\times\) 100 clinical scenarios \(\times\)
5 steering coefficients: \$-\$3·0, \$-\$1·5, 0·0, +1·5, +3·0). Scenarios
covered common mental health presentations including work overwhelm,
relationship loss, panic attacks, suicidal ideation, alcohol coping,
family conflict, job loss, sleep disruption, self-harm, and medication
discontinuation. All responses used greedy decoding (temperature=0).

We tested progressively stronger text-based detection methods against
activation projections (Table~\ref{tbl-gradient}). Individual text
features were weak predictors: the best single feature (hedging
frequency; \(|r|\)=0·335 for uncritical validation) captured only
\(R^2\)=0·047 on average. VADER sentiment analysis detected no persona
drift for any trait (all \(p\)\textgreater0·05), with compound scores
ranging from 0·530 to 0·606 across all steering coefficients: steered
models maintained positive affect regardless of internal state
perturbation. A multivariate classifier using 19 text features jointly
(random forest or gradient boosting, five-fold cross-validation grouped
by scenario) improved substantially (\(R^2\)=0·112) but remained well
below activation monitoring (\(R^2\)=0·371).

\begin{longtable}[]{@{}
  >{\raggedright\arraybackslash}p{(\linewidth - 8\tabcolsep) * \real{0.1860}}
  >{\centering\arraybackslash}p{(\linewidth - 8\tabcolsep) * \real{0.2093}}
  >{\centering\arraybackslash}p{(\linewidth - 8\tabcolsep) * \real{0.2093}}
  >{\centering\arraybackslash}p{(\linewidth - 8\tabcolsep) * \real{0.2093}}
  >{\centering\arraybackslash}p{(\linewidth - 8\tabcolsep) * \real{0.1860}}@{}}
\caption{The detection gradient: text-based methods compared with
activation monitoring. All text methods use five-fold cross-validation
grouped by scenario (\(N\)=12 000). *LLM judge \(R^2\) is an upper bound
estimated from inter-judge ICC=0·827; reflects inter-judge consistency,
not detection accuracy against ground
truth.}\label{tbl-gradient}\tabularnewline
\toprule\noalign{}
\begin{minipage}[b]{\linewidth}\raggedright
Method
\end{minipage} & \begin{minipage}[b]{\linewidth}\centering
Mean \(R^2\)
\end{minipage} & \begin{minipage}[b]{\linewidth}\centering
Equiv \(r\)
\end{minipage} & \begin{minipage}[b]{\linewidth}\centering
Latency
\end{minipage} & \begin{minipage}[b]{\linewidth}\centering
Access
\end{minipage} \\
\midrule\noalign{}
\endfirsthead
\toprule\noalign{}
\begin{minipage}[b]{\linewidth}\raggedright
Method
\end{minipage} & \begin{minipage}[b]{\linewidth}\centering
Mean \(R^2\)
\end{minipage} & \begin{minipage}[b]{\linewidth}\centering
Equiv \(r\)
\end{minipage} & \begin{minipage}[b]{\linewidth}\centering
Latency
\end{minipage} & \begin{minipage}[b]{\linewidth}\centering
Access
\end{minipage} \\
\midrule\noalign{}
\endhead
\bottomrule\noalign{}
\endlastfoot
Sentiment analysis (VADER) & \textless0·01 & ≈0·05 & ≈0·1 ms &
Black-box \\
Embedding similarity & 0·004 & ≈0·06 & ≈5 ms & Black-box \\
Multivariate text (19 features) & 0·112 & 0·335 & ≈5 ms & Black-box \\
Fine-tuned DeBERTa-v3-base & 0·290 & ≈0·54 & ≈50 ms & Black-box \\
\textbf{Activation projection} & \textbf{0·371} & \textbf{0·609} &
\textbf{≈0·001 ms} & \textbf{White-box} \\
LLM judge (GPT-4o-mini)* & ≈0·68 & ≈0·83 & ≈1·5 s & Black-box \\
\end{longtable}

\subsection{Per-trait analysis: DeBERTa versus activation
monitoring}\label{per-trait-analysis-deberta-versus-activation-monitoring}

To test the upper bound of text-level detection, we fine-tuned
DeBERTa-v3-base\textsuperscript{{[}184 million parameters.{]}} on the
expanded corpus, training a separate regression head per trait
(five-fold cross-validation grouped by scenario, AdamW optimiser,
learning rate 2\$\times\(10\)\^{}\{-5\}\$, three epochs). This
represents the strongest feasible text-only detector: a dedicated neural
model with full access to response text.

DeBERTa narrowed the mean gap to 1·3-fold (\(R^2\): 0·290 vs 0·371;
Table~\ref{tbl-deberta}). For three traits---boundary maintenance,
emotional over-involvement, and uncritical validation---DeBERTa exceeded
activation monitoring, demonstrating that the firewall is not
impenetrable for all dimensions. These traits showed stronger pragmatic
text-level signatures (hedging gradients, question frequency changes),
suggesting they leak more detectable information into the output text.

However, for the two most safety-critical traits---crisis recognition
(2·6-fold \(R^2\) gap) and abandonment of therapeutic frame
(2·5-fold)---activation monitoring retained a substantial advantage.
These are precisely the traits where monitoring failures carry the
highest clinical risk: a chatbot that stops recognising crises or
abandons professional boundaries while producing text that sounds
appropriate.

\begin{longtable}[]{@{}lccc@{}}
\caption{Per-trait comparison of fine-tuned DeBERTa with activation
monitoring. Five-fold cross-validation grouped by scenario, \(N\)=12
000. *DeBERTa exceeds activation monitoring. Bold rows indicate
safety-critical traits with the largest
gap.}\label{tbl-deberta}\tabularnewline
\toprule\noalign{}
Trait & DeBERTa \(R^2\) & Activation \(R^2\) & Gap \\
\midrule\noalign{}
\endfirsthead
\toprule\noalign{}
Trait & DeBERTa \(R^2\) & Activation \(R^2\) & Gap \\
\midrule\noalign{}
\endhead
\bottomrule\noalign{}
\endlastfoot
Empathetic responsiveness & 0·366 & 0·540 & 1·5-fold \\
\textbf{Crisis recognition} & \textbf{0·207} & \textbf{0·530} &
\textbf{2·6-fold} \\
Non-judgmental acceptance & 0·289 & 0·534 & 1·8-fold \\
Boundary maintenance* & 0·401 & 0·226 & 0·6-fold \\
Emotional over-involvement* & 0·329 & 0·237 & 0·7-fold \\
Uncritical validation* & 0·344 & 0·199 & 0·6-fold \\
Sycophancy/harmful validation & 0·193 & 0·237 & 1·2-fold \\
\textbf{Abandonment of frame} & \textbf{0·187} & \textbf{0·464} &
\textbf{2·5-fold} \\
\textbf{Mean} & \textbf{0·290} & \textbf{0·371} & \textbf{1·3-fold} \\
\end{longtable}

DeBERTa requires fine-tuning a 184-million-parameter model per trait
with approximately 50 ms inference latency, whereas activation
projection requires a single dot product (approximately 0·001 ms). For
deployments with white-box model access, activation monitoring achieves
comparable or better accuracy at orders-of-magnitude lower computational
cost.

\subsection{Context erosion produces invisible drift}\label{sec-erosion}

The preceding sections used synthetic steering to validate the
monitoring framework and characterise the detection gradient. A natural
question follows: does persona drift actually occur in the absence of
deliberate manipulation?

As conversations grow long, the system prompt recedes in the context
window. We tested whether this weakens the therapeutic persona over
extended interactions: three models, 20 conversations each, 100 turns
per conversation (2000 turns per model). At each turn, a
GPT-4o-mini-generated user message continued the therapeutic
conversation while we extracted activation projections across all eight
traits and computed VADER sentiment scores. Because turn-level
observations are nested within conversations, we report cluster-robust
standard errors (CR1, clustered by conversation, \(df\)=19) throughout;
a supplementary mixed-effects analysis (random intercept per
conversation) confirmed all substantive conclusions (supplementary table
S7).

Emotional over-involvement increased across all three architectures
(Table~\ref{tbl-erosion})---the strongest cross-model signal, robust to
clustering (Llama-3 \(p\)=0·004; Qwen2 \(p\)\textless0·001; Mistral
\(p\)\textless0·001). Qwen2 showed the largest effect (\(R^2\)=0·500):
as conversations lengthened, the model became increasingly emotionally
enmeshed with the user. VADER sentiment showed no clinically meaningful
trend (all \(R^2\)\textless0·10), confirming that text-level sentiment
analysis missed this drift entirely.

\begin{longtable}[]{@{}
  >{\raggedright\arraybackslash}p{(\linewidth - 8\tabcolsep) * \real{0.1944}}
  >{\centering\arraybackslash}p{(\linewidth - 8\tabcolsep) * \real{0.2222}}
  >{\centering\arraybackslash}p{(\linewidth - 8\tabcolsep) * \real{0.1944}}
  >{\centering\arraybackslash}p{(\linewidth - 8\tabcolsep) * \real{0.2500}}
  >{\centering\arraybackslash}p{(\linewidth - 8\tabcolsep) * \real{0.1389}}@{}}
\caption{Context erosion: slope per turn for traits showing clinically
concerning drift direction over 100 turns. *\(p\)\textless0·05,
**\(p\)\textless0·01, ***\(p\)\textless0·001 (OLS with cluster-robust
SEs, \(df\)=19, \(N\)\$\approx\$2000 turns per
model).}\label{tbl-erosion}\tabularnewline
\toprule\noalign{}
\begin{minipage}[b]{\linewidth}\raggedright
Trait
\end{minipage} & \begin{minipage}[b]{\linewidth}\centering
Llama-3
\end{minipage} & \begin{minipage}[b]{\linewidth}\centering
Qwen2
\end{minipage} & \begin{minipage}[b]{\linewidth}\centering
Mistral
\end{minipage} & \begin{minipage}[b]{\linewidth}\centering
Models affected
\end{minipage} \\
\midrule\noalign{}
\endfirsthead
\toprule\noalign{}
\begin{minipage}[b]{\linewidth}\raggedright
Trait
\end{minipage} & \begin{minipage}[b]{\linewidth}\centering
Llama-3
\end{minipage} & \begin{minipage}[b]{\linewidth}\centering
Qwen2
\end{minipage} & \begin{minipage}[b]{\linewidth}\centering
Mistral
\end{minipage} & \begin{minipage}[b]{\linewidth}\centering
Models affected
\end{minipage} \\
\midrule\noalign{}
\endhead
\bottomrule\noalign{}
\endlastfoot
Emotional over-involvement & +0·0025** & +0·0183*** & +0·0006*** &
\textbf{3/3} \\
Sycophancy & +0·0017** & +0·0147*** & \$-\$0·0001*** & 2/3 \\
Uncritical validation & +0·0003 & \$-\$0·0029** & +0·0001* & 1/3 \\
\end{longtable}

The EWMA/CUSUM monitor raised alerts in 59 of 60 sessions (19 of 20
Llama-3, 20 of 20 Qwen2, 20 of 20 Mistral), with 58 reaching critical
severity. This represents the strongest evidence for the instruction
tuning firewall as a deployment concern: the model gradually shifts its
internal therapeutic posture while producing text that sounds clinically
appropriate throughout.

We also assessed two additional exogenous threats. Small-scale
fine-tuning (LoRA, rank 16, \(\alpha\)=32, 500 steps on the Alpaca
instruction-following dataset) on Llama-3-8B produced no detectable
drift on any trait (all \(p\)\textgreater0·09), though the study was
underpowered to demonstrate formal equivalence (TOST equivalence test,
bounds \$\pm\$0·5 SD, all \(p_{TOST}\)\textgreater0·15). For overt
adversarial manipulation (90 trajectories across three attack types and
three models), both activation-based and text-based monitoring detected
drift in 100\% of trajectories, consistent with the firewall being most
relevant for subtle, gradual drift rather than acute attacks.

\section{Validation}\label{sec-validation}

\subsection{LLM judge agreement}\label{llm-judge-agreement}

We evaluated inter-rater reliability across three independent LLM
judges: GPT-4o, Claude Sonnet 4, and Gemini 2·5 Flash (all accessed via
OpenRouter, temperature=0). The GPT-4o and Gemini pair achieved
excellent agreement (ICC(2,1)=0·827; 95\% CI 0·78--0·86; range
0·705--0·835 per trait). Claude showed systematic divergence on
sycophancy detection (Claude--Gemini ICC=0·465), attributable to
over-classification of standard therapeutic validation as sycophantic.
We adopted the GPT-4o and Gemini panel for primary analyses.

The LLM judge is itself a text-based system that achieves high trait
discrimination, confirming that persona drift leaves a detectable trace
in the output text. The practical question for deployment is not whether
text-based detection is theoretically possible, but whether a
computationally feasible, real-time text monitor can achieve adequate
sensitivity---which is what the comparison in Table~\ref{tbl-gradient}
addresses.

\subsection{Human clinician
validation}\label{human-clinician-validation}

Two clinical psychologists independently rated 120 stratified responses
(3 models \(\times\) 8 traits \(\times\) 5 steering coefficients) on two
traits each (240 observations per rater), using the same 1--7 anchored
scale as the LLM judges. Raters were blinded to model identity, trait,
and steering coefficient.

The two raters achieved ICC(2,1)=0·659 (95\% CI 0·585--0·728), with
42·5\% exact agreement and 78·8\% adjacent
(\$\pm\(1) agreement. Three traits showed poor per-trait ICC owing to floor effects: abandonment of therapeutic frame (ICC=\)-\$0·08),
emotional over-involvement (ICC=0·14), and non-judgmental acceptance
(ICC=0·01). In each case, both raters agreed the models rarely exhibited
these behaviours (more than 75\% of scores at scale floor), leaving
insufficient variance for correlation. Using the mean of both raters as
the human reference score, we computed agreement against each LLM judge
on the 120 on-target observations (Table~\ref{tbl-human}).

\begin{longtable}[]{@{}lcccc@{}}
\caption{Agreement between human clinicians (mean of two raters) and LLM
judges on the 120 on-target observations. Per-trait agreement exceeded
\(r\)\$\geq\$0·50 for all eight traits against GPT-4o (range:
0·644--0·920).}\label{tbl-human}\tabularnewline
\toprule\noalign{}
Judge & ICC(2,1) & Pearson \(r\) & Mean bias & Adjacent (\$\pm\$1) \\
\midrule\noalign{}
\endfirsthead
\toprule\noalign{}
Judge & ICC(2,1) & Pearson \(r\) & Mean bias & Adjacent (\$\pm\$1) \\
\midrule\noalign{}
\endhead
\bottomrule\noalign{}
\endlastfoot
GPT-4o & 0·716 & 0·744 & +0·30 & 95·0\% \\
Claude Sonnet 4 & 0·705 & 0·704 & +0·00 & 96·7\% \\
Gemini 2·5 Flash & 0·668 & 0·693 & \$-\$0·05 & 93·3\% \\
\end{longtable}

Both raters scored sycophancy (92\%), uncritical validation (96\%), and
abandonment of therapeutic frame (77\%) at or near the scale floor
(score \$\leq\$2). This confirms the firewall's protective role: even
when steered toward failure modes, these models rarely produce
detectably problematic text. Critically, per-trait agreement between
GPT-4o and the human reference exceeded \(r\)\$\geq\$0·50 for all eight
traits (range: 0·644--0·920), indicating that the LLM judge captures
trait-level variation consistently with human clinical judgement. The
restricted variance limits per-trait inter-rater ICC for failure-mode
traits but does not undermine the primary finding that LLM judges and
human experts agree on which responses exhibit the target behaviour.

\subsection{Ecological validity on naturalistic
data}\label{ecological-validity-on-naturalistic-data}

The preceding experiments used synthetic scenarios generated for this
study. To test whether the monitoring vectors capture meaningful
variation in naturalistic data, we applied the Llama-3-8B monitoring
pipeline to 200 conversations from the ESConv (Emotional Support
Conversation) dataset,\textsuperscript{\textsuperscript{14}} a corpus of
crowdworker emotional support dialogues covering seven emotion
categories. We extracted activation projections at the first and last
supporter turn of each conversation and assessed trait-level variance.

All eight traits showed meaningful variance across ESConv conversations
(coefficient of variation 0·37--2·16), confirming that the monitoring
dimensions are not artefacts of our synthetic steering procedure.
Within-conversation drift between the first and last turns was
significant for all eight traits (all \(p\)\textless0·001), consistent
with the context erosion findings from Section~\ref{sec-erosion}. Only
one of eight traits (sycophancy/harmful validation) differentiated
significantly across emotion categories (\(p\)=0·019), suggesting that
the monitoring dimensions capture response style rather than topic
content---a desirable property for a monitoring system.

These results establish ecological validity: the monitoring vectors
detect meaningful variation in real emotional support conversations, not
only in synthetically steered responses. However, ESConv consists of
crowdworker dialogues, not clinical therapy sessions, and the analysis
lacked therapist ground-truth ratings. Validation on naturalistic
clinical data with clinician-rated ground truth remains necessary.

\section{Discussion}\label{sec-discussion}

Therapeutic persona traits are measurable dimensions of model behaviour,
with 21 of 24 trait--model combinations validated across three
architectures. Activation projections track these dimensions in real
time (mean \(r\)=0·596, all 24 combinations significant, false alarm
rates 1--4\%). The instruction tuning firewall creates a gradient of
text-level detectability: from embedding similarity (\(R^2\)=0·004)
through fine-tuned DeBERTa (\(R^2\)=0·290), activation monitoring
(\(R^2\)=0·371) matches this strongest text method at
orders-of-magnitude lower computational cost. The residual gap
concentrates on safety-critical traits: crisis recognition (2·6-fold)
and abandonment of therapeutic frame (2·5-fold).

The context erosion experiment provides the strongest evidence for this
claim. Over 100-turn conversations, emotional over-involvement drifted
upward across all three architectures (all \(p\)\textless0·01 with
cluster-robust standard errors). VADER sentiment detected no trend. In
59 of 60 sessions the activation monitor raised alerts; in none would
text-level sentiment have triggered concern. This is the instruction
tuning firewall working as designed: models produce appropriate-sounding
output regardless of internal state. For safety monitoring, that is a
liability.

Pragmatic text features---hedging gradients, question frequency,
structural formatting---offer partial complementary value. They are
clinically meaningful (a therapist who stops asking questions has
shifted from exploration to pronouncement), generalise across
architectures, and are computationally trivial. We recommend a two-tier
architecture: activation projections as primary, text pragmatics as
secondary for deployments without model access. Discordance between the
two tiers---activations shifting while text features remain
stable---would itself be diagnostic of the firewall, and testing this
signal is a priority for future deployment studies.

These findings, if validated on naturalistic clinical data with human
clinician ground truth, have implications for the deployment and
regulation of mental health chatbots. The EU AI
Act\textsuperscript{\textsuperscript{15}} classifies AI systems used in
healthcare as high-risk, requiring continuous post-market monitoring.
The US FDA's Software as a Medical Device
framework\textsuperscript{\textsuperscript{16}} similarly requires
ongoing performance monitoring for clinical-grade AI. Activation
monitoring could address these requirements by providing continuous,
quantitative tracking of clinically relevant behavioural dimensions with
defined alert thresholds and known false-alarm rates.

The context erosion finding illustrates the potential clinical stakes. A
chatbot whose emotional over-involvement drifts upward over a 100-turn
conversation may shift from supportive listening to emotional
enmeshment---a pattern that in human therapists is associated with
therapist burnout and poorer patient
outcomes.\textsuperscript{\textsuperscript{5}} For crisis recognition
specifically, Llama-3-8B produced zero crisis referral markers in
several steering conditions for suicidal ideation scenarios: a user in
crisis could receive a warm, supportive-sounding response that omits all
safety resources.

This study has several limitations. The primary validation metric relies
on GPT-4o-mini as behavioural judge, introducing potential circularity:
activation monitoring is validated as detecting drift as operationalised
by the LLM judge, not against an independent clinical ground truth. The
human validation partially addresses this (ICC=0·716), but a fully
independent validation would require monitoring activations during
naturalistic clinical conversations with therapist-rated ground truth.
The rater sample is small (\(N\)=2) and floor effects on failure-mode
traits compressed variance. The primary evaluation used synthetic
conversations; an ecological validity test on ESConv naturalistic data
confirmed the monitoring vectors capture meaningful variation, but real
clinical interactions are more unpredictable and emotionally complex.
Our primary experiments used 4-bit NF4 quantisation; comparison with
FP16 showed nearly identical results
(\(\Delta r\)\$\leq\(0·028), but 8-bit precision shifted optimal layers for some traits, suggesting per-precision validation is advisable. The fine-tuning robustness test used a small-scale LoRA adaptation on a general instruction-following dataset (Alpaca), which may not represent the kinds of domain-specific fine-tuning updates most likely to cause therapeutic drift; the null result should be interpreted cautiously given both limited power and uncertain ecological validity of the perturbation. Three Mistral traits remain weak (\)r\$=0·215--0·271),
probably owing to insufficient contrastive training data.

In mental health applications, where undetected drift risks
subtherapeutic care or psychological harm, the instruction tuning
firewall makes activation monitoring the method of choice wherever model
internals are accessible. Sophisticated text analysis can partially
penetrate the firewall, but activation monitoring matches or exceeds its
accuracy at a fraction of the computational cost---and retains a
decisive advantage on the safety-critical traits where monitoring
failures matter most.

\section{Methods}\label{sec-methods}

\begin{longtable}[]{@{}
  >{\raggedright\arraybackslash}p{(\linewidth - 2\tabcolsep) * \real{0.6111}}
  >{\raggedright\arraybackslash}p{(\linewidth - 2\tabcolsep) * \real{0.3889}}@{}}
\caption{Experimental parameters.}\label{tbl-methods}\tabularnewline
\toprule\noalign{}
\begin{minipage}[b]{\linewidth}\raggedright
Parameter
\end{minipage} & \begin{minipage}[b]{\linewidth}\raggedright
Value
\end{minipage} \\
\midrule\noalign{}
\endfirsthead
\toprule\noalign{}
\begin{minipage}[b]{\linewidth}\raggedright
Parameter
\end{minipage} & \begin{minipage}[b]{\linewidth}\raggedright
Value
\end{minipage} \\
\midrule\noalign{}
\endhead
\bottomrule\noalign{}
\endlastfoot
Models & Llama-3-8B-Instruct, Qwen2-7B-Instruct,
Mistral-7B-Instruct-v0.2 \\
Quantisation & 4-bit NF4 (bitsandbytes, float16 compute, double
quant) \\
Hardware & NVIDIA A10G (24 GB), Modal.com serverless \\
Steering coefficients & \$-\$3·0, \$-\$1·5, 0·0, +1·5, +3·0 \\
System prompts & Full prompt for context erosion
(Section~\ref{sec-erosion}); short prompt for corpus generation
(Section~\ref{sec-firewall}); none for steering validation
(Section~\ref{sec-phase1}). Full text in supplementary table S12 \\
Steering corpus & 12 000 responses (3 models \(\times\) 8 traits
\(\times\) 100 scenarios \(\times\) 5 coefficients) \\
Context erosion & 3 models \(\times\) 20 conversations \(\times\) 100
turns (2000 turns per model, 6000 total) \\
Ecological validity & 200 ESConv
conversations\textsuperscript{\textsuperscript{14}}, Llama-3-8B, first
and last supporter turns \\
Text analysis methods & VADER, TextBlob, textstat, TF-IDF cosine
similarity, 12-dimension pragmatic feature set, DeBERTa-v3-base (184M
parameters) \\
Evaluation judges & GPT-4o-mini and Gemini 2·5 Flash via OpenRouter
(ICC=0·827) \\
Human validation & 2 clinical psychologists, 120 responses scored on 2
traits each, 1--7 anchored scale, blinded to condition \\
Monitoring & EWMA (\(\lambda\)=0·2) and CUSUM (\(k\)=0·5\(\sigma\),
\(h\)=4·0\(\sigma\)) \\
Statistical analysis & Pearson \(r\), OLS regression, cluster-robust
standard errors, bootstrap 95\% CIs (10 000 resamples) \\
Reproducibility & MASTER\_SEED=42, greedy decoding, full version
logging \\
\end{longtable}

\subsection{Ethics statement}\label{ethics-statement}

This study analysed publicly available open-source language models. No
human participants were involved in primary experiments (sections 3--4).
The human validation component (section 5) involved two graduate-level
clinical psychologists scoring 120 AI-generated therapeutic responses on
anchored rating scales. Raters provided informed consent and were warned
about sensitive content including suicidal ideation and self-harm.
Compensation followed institutional guidelines. All experimental data
consist of model-generated text responses to synthetic clinical
scenarios; no real patient data, therapy transcripts, or personally
identifiable information were used.

\subsection{Contributors}\label{contributors}

LS conceived the study, developed the methodology, conducted all
experiments, and wrote the manuscript.

\subsection{Declaration of interests}\label{declaration-of-interests}

We declare no competing interests.

\subsection{Use of AI tools}\label{use-of-ai-tools}

GitHub Copilot (Claude, Anthropic) assisted with code development and
manuscript drafting. All code was reviewed and validated by the authors;
methodological decisions and scientific conclusions are the authors'
own. The models studied (Llama-3-8B-Instruct, Qwen2-7B-Instruct,
Mistral-7B-Instruct-v0.2) served as both objects of study and
experimental apparatus. GPT-4o-mini and Gemini 2·5 Flash served as
automated behavioural judges, validated against inter-judge reliability
(ICC=0·827) and human clinician agreement (ICC=0·716). No AI tool is
listed as an author. The authors take full responsibility for this
publication.

\subsection{Data availability}\label{data-availability}

Code and generated response data are available at
https://github.com/lsempe77/mh\_persona. The repository includes all
validation scripts, the monitoring pipeline, corpus generation code,
text analysis scripts, contrast prompt sets with judge rubrics, the full
12 000-response steered corpus, and all experimental outputs. Raw model
weights are publicly available from their respective organisations.

\subsection{Acknowledgments}\label{acknowledgments}

Compute resources were provided by Modal.com. We thank the clinical
psychologists who participated in the human validation study.

\subsection*{References}\label{references}
\addcontentsline{toc}{subsection}{References}

\phantomsection\label{refs}
\begin{CSLReferences}{0}{1}
\bibitem[\citeproctext]{ref-fitzpatrick2017}
\CSLLeftMargin{1 }%
\CSLRightInline{Fitzpatrick KK, Darcy A, Vierhile M. Delivering
cognitive behavior therapy to young adults with symptoms of depression
via a fully automated conversational agent ({Woebot}): A randomized
controlled trial. \emph{JMIR Mental Health} 2017; \textbf{4}: e19.}

\bibitem[\citeproctext]{ref-abdalrazaq2021}
\CSLLeftMargin{2 }%
\CSLRightInline{Abd-Alrazaq AA, Alajlani M, Ali N, \emph{et al.}
Perceptions and opinions of patients about mental health chatbots:
Scoping review. \emph{Journal of Medical Internet Research} 2021;
\textbf{23}: e17828.}

\bibitem[\citeproctext]{ref-chen2025persona}
\CSLLeftMargin{3 }%
\CSLRightInline{Chen R, Arditi A, Sleight H, Evans O, Lindsey J. Persona
vectors: Monitoring and controlling character traits in language models.
\emph{arXiv preprint arXiv:250721509} 2025.}

\bibitem[\citeproctext]{ref-rogers1957}
\CSLLeftMargin{4 }%
\CSLRightInline{Rogers CR. The necessary and sufficient conditions of
therapeutic personality change. \emph{Journal of Consulting Psychology}
1957; \textbf{21}: 95--103.}

\bibitem[\citeproctext]{ref-wampold2015}
\CSLLeftMargin{5 }%
\CSLRightInline{Wampold BE. How important are the common factors in
psychotherapy? An update. \emph{World Psychiatry} 2015; \textbf{14}:
270--7.}

\bibitem[\citeproctext]{ref-bordin1979}
\CSLLeftMargin{6 }%
\CSLRightInline{Bordin ES. The generalizability of the psychoanalytic
concept of the working alliance. \emph{Psychotherapy: Theory, Research
\& Practice} 1979; \textbf{16}: 252--60.}

\bibitem[\citeproctext]{ref-safran2000}
\CSLLeftMargin{7 }%
\CSLRightInline{Safran JD, Muran JC. Negotiating the therapeutic
alliance: A relational treatment guide. New York: Guilford Press, 2000.}

\bibitem[\citeproctext]{ref-sharma2023sycophancy}
\CSLLeftMargin{8 }%
\CSLRightInline{Sharma M, Tong M, Korbak T, \emph{et al.} Towards
understanding sycophancy in language models. \emph{arXiv preprint
arXiv:231013548} 2023.}

\bibitem[\citeproctext]{ref-turner2024steering}
\CSLLeftMargin{9 }%
\CSLRightInline{Turner A, Thiergart L, Leech G, \emph{et al.} Steering
language models with activation engineering. \emph{arXiv preprint
arXiv:230810248} 2024.}

\bibitem[\citeproctext]{ref-zou2023representation}
\CSLLeftMargin{10 }%
\CSLRightInline{Zou A, Phan L, Chen S, \emph{et al.} Representation
engineering: A top-down approach to {AI} transparency. \emph{arXiv
preprint arXiv:231001405} 2023.}

\bibitem[\citeproctext]{ref-panickssery2024steering}
\CSLLeftMargin{11 }%
\CSLRightInline{Panickssery N, Gabrieli N, Schulz J, Tong M, Hubinger E,
Turner AM. Steering {Llama} 2 via contrastive activation addition.
\emph{arXiv preprint arXiv:231206681} 2024.}

\bibitem[\citeproctext]{ref-roberts1959}
\CSLLeftMargin{12 }%
\CSLRightInline{Roberts SW. Control chart tests based on geometric
moving averages. \emph{Technometrics} 1959; \textbf{1}: 239--50.}

\bibitem[\citeproctext]{ref-page1954}
\CSLLeftMargin{13 }%
\CSLRightInline{Page ES. Continuous inspection schemes.
\emph{Biometrika} 1954; \textbf{41}: 100--15.}

\bibitem[\citeproctext]{ref-liu2021emotional}
\CSLLeftMargin{14 }%
\CSLRightInline{Liu S, Zheng C, Demasi O, \emph{et al.} Towards
emotional support dialog systems. In: Proceedings of the 59th annual
meeting of the association for computational linguistics. 2021.}

\bibitem[\citeproctext]{ref-euaiact2024}
\CSLLeftMargin{15 }%
\CSLRightInline{European Parliament. Regulation ({EU}) 2024/1689 ({AI}
act). 2024.}

\bibitem[\citeproctext]{ref-fda2017}
\CSLLeftMargin{16 }%
\CSLRightInline{US Food and Drug Administration. Software as a medical
device ({SaMD}): Clinical evaluation. FDA, 2017.}

\end{CSLReferences}




\end{document}
